\documentclass{standalone}

\usepackage{tikz}

%% Define custom color palatte
\definecolor{primary}{HTML}{dddddd}
\definecolor{secondary}{HTML}{aaaaaa}
\definecolor{tertiary}{HTML}{666666}

%% Define constants & functions
\def\r{5}
\newcommand{\pathtwo}[2]{
    (d#1) to [out=270-#1*36, in=270-#2*36] (d#2)
}
\newcommand{\paththree}[3]{
    (d#1) to [out=270-#1*36, in=270-#2*36] (d#2) to [out=270-#2*36, in=270-#3*36] (d#3)
}
\newcommand{\fillbg}[1]{
    \fill[#1] (\r,\r) rectangle (-\r,-\r);
}

\begin{document}
\begin{tikzpicture}

\pgfmathparse{tan(36)}
\coordinate (d0) at (0, \r);
\coordinate (d5) at (0,-\r);
\coordinate (d1) at ( \pgfmathresult*\r, \r);
\coordinate (d9) at (-\pgfmathresult*\r, \r);
\coordinate (d4) at ( \pgfmathresult*\r,-\r);
\coordinate (d6) at (-\pgfmathresult*\r,-\r);

\pgfmathparse{tan(18)}
\coordinate (d2) at ( \r, \pgfmathresult*\r);
\coordinate (d3) at ( \r,-\pgfmathresult*\r);
\coordinate (d7) at (-\r,-\pgfmathresult*\r);
\coordinate (d8) at (-\r, \pgfmathresult*\r);

%% pi-10 Wireframe
% \def\pivalue{{3,1,4,1,5,9,2,6,5,3,5,8}}
% \foreach \i in {0,...,10}
% {
%     \pgfmathsetmacro{\j}{\i+1}
%     \pgfmathsetmacro{\ipi}{\pivalue[\i]}
%     \pgfmathsetmacro{\jpi}{\pivalue[\j]}
%     \draw (d\ipi) to [out=270-\ipi*36,in=270-\jpi*36] (d\jpi);
% }

\begin{scope}
    \clip \pathtwo{9}{2} -- (\r,\r) -- cycle;
    \fillbg{primary};
    \begin{scope}
        \clip \paththree{4}{1}{5} -- cycle;
        \fillbg{tertiary};
    \end{scope}
    \begin{scope}
        \clip \paththree{3}{1}{4} -- cycle;
        \fillbg{secondary};
    \end{scope}
\end{scope}
\begin{scope}
    \clip \pathtwo{2}{9} -- (-\r,\r) -- (-\r,-\r) -- \pathtwo{6}{2} -- cycle;
    \fillbg{tertiary};
    \begin{scope}
        \clip \pathtwo{1}{3} -- (\r,\r) -- cycle;
        \fillbg{secondary};
    \end{scope}
    \begin{scope}
        \clip \paththree{4}{1}{5} -- cycle;
        \fillbg{secondary};
    \end{scope}
    \begin{scope}
        \clip \paththree{8}{5}{9} -- (-\r,\r) -- cycle;
        \fillbg{secondary};
    \end{scope}
\end{scope}
\begin{scope}
    \clip \pathtwo{2}{6} -- \pathtwo{5}{3} -- cycle;
    \fillbg{secondary};
    \begin{scope}
        \clip \pathtwo{1}{3} -- (\r,\r) -- cycle;
        \fillbg{tertiary};
    \end{scope}
    \begin{scope}
        \clip \paththree{4}{1}{5} -- cycle;
        \fillbg{primary};
    \end{scope}
    \begin{scope}
        \clip \paththree{8}{5}{9} -- cycle;
        \fillbg{tertiary};
    \end{scope}
    \begin{scope}
        \clip \pathtwo{5}{6} -- cycle;
        \fillbg{primary};
    \end{scope}
\end{scope}
\begin{scope}
    \clip \pathtwo{3}{5} -- (\r,-\r) -- cycle;
    \fillbg{primary};
    \begin{scope}
        \clip \paththree{4}{1}{5} -- cycle;
        \fillbg{tertiary};
    \end{scope}
\end{scope}
\end{tikzpicture}
\end{document}